\section{有限温度微扰论}

Gross 第 29 章把零温图技术推广到有限 $T$：用密度矩阵/配分函数代替真空期望，Matsubara 频率自然出现。

\subsection{从配分函数出发}

\begin{equation}
  Z
  =\mathrm{Tr}\,e^{-\beta(H_0+V)}
  =Z_0
  \Bigl\langle
  T_\tau\exp\Bigl(
  -\int_0^\beta\mathrm{d}\tau\,V_I(\tau)
  \Bigr)
  \Bigr\rangle_0.
\end{equation}
展开后 Wick 收缩给出虚时图；对时间积分等价于 Matsubara 频率守恒。连通图指数化后，$\Omega-\Omega_0=-\beta^{-1}\ln(Z/Z_0)$ 等于连通真空图之和。

\subsection{与零温的对应}

$\beta\to\infty$ 且化学势固定在能隙内时，Matsubara 和变为实频积分，连锁集群回到 Goldstone。有限 $T$ 的优势是：无需绝热开关即可定义热力学势，并自然包含热激发（费米/玻色分布）。实践中第 3 章 Matsubara 规则已覆盖计算层面；此处强调概念来源是配分函数而非真空振幅。
